\documentclass[titlepage,a4paper]{jsarticle}
\usepackage{../../../sty/import}% 各種パッケージインポート
\usepackage{../../../sty/title}% タイトルページの変更
\usepackage{../../../sty/answergrid}
%% タイトルページの変数
% レポートタイトル
\title{心理学特論第8回目出席票}
% 提出日
\expdate{\today}
% 科目名
\subject{心理学特論}
% 分野
\class{情報経営システム工学分野}
% 学年
\grade{M1}
% 学籍番号
\mynumber{24336488}
% 記述者
\author{本間三暉}

\begin{document}
% titleページ作成
\maketitle
\section*{keyword\#1}
コーピング
\section*{keyword\#2}
ストレスチェック
\section*{心理学特論第08回の出席課題は，心理学特論サイコロワークと連動しています。2つの事例問題を読んでから取り組んで下さい。}
\section{ケース01甲さんの事例について，メンタルヘルスの観点から，あなたの意見を書いて下さい。（自分ならばこのようにしたかも知れないといったものでも構いません）}
ケース01について，甲さんは上司の異動によって業務内容や職場環境が大きく変化し，多くのストレスを抱えていたと考えられる．
特に，以前は上司が担当していたクレーム対応や他部署との調整を自分一人で担うようになったことで，精神的負担が増加したと思われる．
その結果，胃痛や生理不順といった身体的症状だけでなく，イライラや無気力，仕事に行きたくないという心理面の変化も現れていた．
これは授業で扱われたストレス反応やメンタルヘルス不調の例に当てはまると考えられる．

また，甲さんは環境の変化に対して受け身で耐え続けていたように見える．
自分ならば，一人で我慢を続けるのではなく，上司に業務量について相談したり，周囲に助けを求めたりして，少しでも自分が働きやすい環境になるように行動したと思う．
授業資料でも，問題そのものを改善しようとする「問題焦点型コーピング」が重要であると説明されており，ストレス要因を減らすための行動が必要であると感じた．

さらに，契約社員になった後に仕事をほとんど任されなくなったことで，甲さんは自分の役割や存在意義を感じにくくなり，より無気力になってしまったのではないかと思う．
そのため，退職を考える前に，信頼できる人や専門家に相談し，自分のストレス状態を客観的に把握することが大切だったのではないかと考える．

\section{ケース02乙さんの事例について，メンタルヘルスの観点から，あなたの意見を書いて下さい。（自分ならばこのようにしたかも知れないといったものでも構いません）}
ケース02について，乙さんは結婚や仕事の繁忙期，上司との人間関係の悪化，さらに父親の入院など，多くのストレス要因が重なっていたと考えられる．
特に，夫である丙さんも仕事で大きな問題を抱えており，お互いに精神的余裕が少ない状態だったことが，さらに負担を大きくしていたと思われる．

また，乙さんは「相手のほうが大変だから」と考えて，自分の悩みを相談できずに一人で抱え込んでいた．
しかし，自分は「相手が大変だからといって，自分が大変ではないことにはならない」と考えている．
そのため，自分ならば，お互いに仕事の悩みやつらさを共有し，「お前も大変かもしれないけど，俺も大変なんだ」という形で支え合おうとしたと思う．
授業資料でも，助けを求めることや周囲のサポートを活用することの重要性が説明されており，一人で抱え込まないことが大切だと感じた．

さらに，乙さんには頭痛や肩こりといった身体面の変化だけでなく，仕事でのミスや作業効率の低下といった行動面の変化も現れていた．
これはストレスによる影響が身体や行動にまで広がっている状態であり，かなり強いストレス状態だったと考えられる．

そのため，無理を続けるのではなく，休息や気分転換の時間を確保しながら，周囲と悩みを共有することが重要だったのではないかと思う．
また，自分一人だけで耐えようとするのではなく，互いに支え合う関係を作ることが，メンタルヘルスの維持につながると感じた．

\newpage
\section{心理学特論課題サイコロワーク解答記入欄}
\AnswerGridFilled{
  {確証バイアス}
  {固着性ヒューリスティック}
  {フレーミング効果}
  {アンカリング効果}
  {ハロー効果}
  {セリエ}
  {ストレッサー}
  {汎適応症候群}
  {コーピング}
  {問題焦点}
  {エ}
  {イ}
  {ウ}
  {イ}
  {エ}
  {イ}
  {ア}
  {ウ}
  {ア}
  {イ}
}

\end{document}